% Folder 84, batch 02 (archivists' pages 21-40).
% Pass: Opus 5.5 (claude-opus-5-5), 2026-10-09 — first pass, unchecked against the pages by a human.
%
\documentclass[11pt,a4paper]{article}
\input{../preamble/grothendieck}

\folder{84}
\batch{2}
\pages{21}{40}
\foldertitle{[Formes quadratiques] : notes manuscrites (s.d.).}
\dating{[à partir de 1982-vers 1986]}
\watermark{Édition de démonstration}

\begin{document}

\page{21}
\note{la page s'ouvre sur un N.B. dont les points se rapportent à la construction de l'algèbre $A_Q$ attachée à un triple $(E, \Delta, Q)$, faite avant le début du lot}
\emph{NB} 1°. La corr.\ $(E, \Delta, Q) \mapsto \{A_Q,$ \struck{$A_Q \to \underline{\mathrm{End}}(E)$}\uncertain{$\}$} \add{(pour les iso)} \ill{} \uncertain{famille}, et l'hom.\ \uncertain{(pour mon.\ inj.)}
\[
  A_Q \to \underline{\mathrm{End}}(E)
\]
est fonctoriel. De plus, ces constructions \uncertain{commutent} aux changements de base.

2°. Il en résulte, par transport de structure, \uncertain{qu'à tout} $\lambda \in \Gamma(\mathcal{O}_X^{*})$ (section \emph{inversible}), on \ill{} un isom.\ can.\ d'alg.\ quadratiques
\[
  A_Q \simeq A_{\lambda Q}
\]
compatible avec les actions sur $E$. Bien sûr, il n'en est pas ainsi si on ne suppose pas $\lambda$ inversible (prendre $\lambda = 0$ !).

3°) Conditions équivalentes :
\begin{enumerate}
  \item[a)] $A_Q \to \underline{\mathrm{End}}(E)$ univ.\ inj.\ (i.e.\ inj.\ sur chaque fibre) ;
  \item[b)] $Q$ non nulle sur aucune fibre \uncertain{résiduelle} ;
  \item[c)] $\Delta \xrightarrow{Q} \underline{\mathrm{Quad}}(E)$ est univ.\ injectif ;
  \item[d)] $\Sigma Q$ est structure fidèle.
\end{enumerate}
(C'est essentiellement le lemme 2 plus haut.) Alors la structure d'Alg.\ sur $A_Q$ est \uncertain{simplement} celle induite par $\underline{\mathrm{End}}(E)$.

2.5. Cas \uncertain{tot.} dégénéré de la forme \emph{nulle} $Q \colon E \to \Delta'$. Alors
\[
  A_Q = \underline{\mathcal{O}} \times \Lambda \qquad (\Lambda = \Delta \otimes L)
\]
avec la structure d'algèbre définie par la \uncertain{condition} \uncertain{que} $\Lambda$ est de \emph{carré nul}, i.e.

\page{22}
\[
  A_Q \simeq \mathrm{D}\ill{}^{*}_{\underline{\mathcal{O}}}(\Lambda) .
\]
\note{le nom de l'algèbre n'est pas lu ; d'après la phrase précédente il s'agit de l'algèbre des nombres duaux sur $\Lambda$, $\underline{\mathcal{O}} \oplus \Lambda$ avec $\Lambda^{2} = 0$}
L'hom.
\[
  u_Q^{!} \colon A_Q \to \underline{\mathrm{End}}(E)
\]
est \struck{défini} \add{déterminé} par la condition \struck{que} de s'annuler sur $\Lambda$, i.e.\ c'est $(a, \lambda) \mapsto a$. \struck{\ill{}} Il n'a rien d'injectif.

2.6. Quand on a donné une base de $E$ et de $\Delta$, donc $Q$ est donnée par
\[
  Q(\lambda, \mu) = \alpha \lambda^{2} + \beta \lambda \mu + \gamma \mu^{2}, \qquad \alpha, \beta, \gamma \in \Gamma(\underline{\mathcal{O}}),
\]
et on \struck{peut prendre} a une \struck{base} $(1, u)$ de $A_Q$, où
\[
  \left|\begin{array}{l}
    u^{2} - \beta u + \alpha\gamma = 0 \\
    \mathrm{Tr}\, u = \beta \\
    \mathrm{N}\, u = \alpha\gamma
  \end{array}\right.
\]
i.e.\ \struck{$\beta = \mathrm{Tr}\, u$}. L'image de $u$ dans $\underline{\mathrm{End}}(E) \simeq \mathbb{M}_2(\underline{\mathcal{O}})$ est la matrice déjà envisagée
\[
  u_Q = \begin{pmatrix} 0 & -\gamma \\ \alpha & \beta \end{pmatrix} .
\]

Pour achever de déterminer $A_Q$ comme foncteur en $(E, \Delta, Q)$, il faut \uncertain{encore} expliciter l'isom.\ de
\[
  A_{\alpha, \beta, \gamma} = A[u]/(u^{2} - \beta u + \alpha\gamma)
\]
avec
\[
  A_{\alpha', \beta', \gamma'} = A[u]/(u^{2} - \beta' u + \alpha'\gamma'),
\]
quand on passe d'un syst.\ de bases de $(A, \Delta)$\note{sic, pour $(E, \Delta)$ sans doute} à un autre, par des matrices $\begin{pmatrix} u_{11} & u_{12} \\ u_{21} & u_{22} \end{pmatrix}$ et $v$…

\page{23}
Je ne vais pas faire l'explicitation ici.

2.6. L'exten.\ quadratique \add{de $X$} associée à $A_Q$ mérite de s'appeler \emph{l'invariant d'Arf} de la forme quadratique binaire $Q$. On définit un inv.\ d'Arf pour \emph{les formes} quadratiques scalaires sur un module loc.\ libre de rang donné (mais en distinguant les cas du rang pair et du rang impair, pour la notion de cet invariant…).\note{deux paragraphes portent le numéro 2.6, p.~22 et ici}

\section*{3. Théorème récapitulatif}

Équivalences entre \struck{\ill{}} catégories fibrées sur \add{la catégorie} \struck{Sch} des schémas $X$ :

① Structures de similitude \struck{binaires} \add{(str.\ fidèles)} sur $X$, i.e.\ des systèmes
\[
  (E, \Delta),
\]
où $E$ est un \struck{Module} loc.\ libre de rang 2 sur $X$, et
\[
  \Delta \subset \underline{\mathrm{Quad}}(E)
\]
un sous-\struck{Module} \add{fibré} \struck{vectoriel} \uncertain{vectoriel} inversible de $\underline{\mathrm{Quad}}(E)$, i.e.\ un sous-module inversible et loc.\ facteur direct (ou encore, tel que $\Delta \to \underline{\mathrm{Quad}}\, E$ soit univ.\ injectif).

\page{24}
①\('\) Systèmes
\[
  (E, \Delta', Q),
\]
où $E$ est comme dessus, $\Delta'$ un Module inv.,
\[
  Q \colon E \to \Delta'
\]
une forme quadratique, telle que $Q$ soit non \add{(identiquement)} nulle sur \struck{\ill{}} \uncertain{toute} fibre résiduelle.

② Systèmes \struck{\ill{}} $(A, E)$ d'une alg.\ quadratique $A$ sur $\underline{\mathcal{O}}$, et d'un $A$-module inversible i.e.\ loc.\ libre de rang 1.

②\('\) Systèmes $(X', \Lambda')$ d'un rev.\ quadratique $X'$ de $X$, et d'un Module inversible $\Lambda'$ sur $X'$.

\emph{NB} $1 \Leftrightarrow 1'$ trivial, $2 \Leftrightarrow 2'$ est standard par dictionnaire des spectres d'Anneaux \uncertain{quasi-coh.} $(1, 1') \Rightarrow 2$ a été vu plus haut. Il faut expliciter ②\('\) $\to$ ①\('\). On prend
\[
  \left\{\begin{array}{l}
    E = p_{*}(\Lambda'), \quad \text{où } p \colon X' \to X \text{ morph.\ struct.} \\
    \Delta' = \mathrm{N}_{X'/X}\, \Lambda' \\
    Q \colon p_{*}(\Lambda') \to \Delta' = \mathrm{N}_{X'/X}(\Lambda')
  \end{array}\right.
\]
l'\struck{\ill{}}application norme.

\page{25}
\note{« 30 » cerclé en haut à droite, sans doute la pagination de l'auteur}
\emph{Cor.\ 1} Ces catégories équivalent \uncertain{aussi} à celle des systèmes
\[
  (A, P),
\]
où $A$ est une alg.\ quadratique sur $X$, et où $P$ est un torseur* sous le schéma en groupes \uncertain{lisse} sur $S$ \add{de dimension relative 2}
\[
  W(A)^{*} = \{ S/X \mapsto \Gamma(S, A_S^{*}) \}
\]
(où $A_S$ désigne l'image inverse de $A$ sur $S$).
\marginal{* torseur pour la \uncertain{top.\ fpqc}, \uncertain{ou} pour (Zar), et cela revient \uncertain{au même} \ill{}, \uncertain{pour toute topologie intermédiaire} \ill{} \uncertain{schémas}.}

En termes des $(E, Q)$, $P$ est le \struck{En termes des} schéma sur $X$ des sections locales de $E$, telles que $Q(x)$ soit une base de $\Delta'$. D'autre part, $\Delta'$ s'explicite en termes de $(A, P)$ comme le fibré en droites \uncertain{déduit} du fibré principal associé
\[
  P \overset{W(A)^{*}}{\wedge} \mathbb{G}_m ,
\]
via l'hom.\ norme
\[
  W(A)^{*} \xrightarrow{\mathrm{N}_{A/\underline{\mathcal{O}}}} \mathbb{G}_m ,
\]
qui s'insère dans la suite exacte
\[
  1 \to \mathcal{U}_A \to W(A)^{*} \xrightarrow{\mathrm{N}} \mathbb{G}_m \to 1,
\]
$\mathcal{U}_A$ = schéma en groupes « unitaire » des éléments de norme 1 de $A$.

\page{26}
\emph{Cor.\ 2} Équiv.\ de cat.\ fibrées entre

1°) $(E, Q)$, $Q$ une forme quadratique binaire qui n'est nulle sur aucune fibre résid.

2°) Systèmes $(A, \Pi)$, où $A$ est une alg.\ quadratique sur $S$, et où $\Pi$ est un torseur* sous le groupe unitaire $\mathcal{U}_A$.
\marginal{* Torseur \uncertain{sous} \ill{} fpqc ou \uncertain{sous} (ét), cela revient au même, mais la top.\ (Zar) n'est plus \uncertain{assez fine}.}

$\Pi$ est le sous-schéma de $\mathbb{E}$\note{un $E$ surmonté d'un signe, sans doute le fibré vectoriel associé à $E$} des sections locales $x$ telles que $Q(x) = 1$, qui \struck{forment} \add{est bien} un torseur sous $\mathcal{U}_A$.

\emph{NB} Le rev.\ quadratique associé à la forme quadratique $Q \colon E \to \Delta'$ (non nulle sur toute fibre) peut s'expliciter comme \add{le} sous-schéma de $\mathbb{P}^{1}(E)$ défini par l'équation $Q(x) = 0$. Si $E$ est considéré comme Module inv.\ sur l'Algèbre quadratique $A$, le morphisme
\[
  \mathrm{Spec}(A) \hookrightarrow \mathbb{P}^{1}(E)
\]
correspondant \struck{\ill{}} s'explicite ainsi : pour toute section locale $x$ de $\mathrm{Spec}(A)$ sur $A$ i.e.

\page{27}
pour tout hom.\ d'algèbres $A \to \underline{\mathcal{O}}$, le noyau $J \subset A$ est un sous-fibré vectoriel \add{de rang 1} \struck{(\ill{} suppl.\ de $\underline{\mathcal{O}} \cdot 1_A$)} \add{\uncertain{(cet inconvénient)}}\note{lecture très douteuse de cet ajout interlinéaire} \add{sur lequel $\mathrm{N}_{A/\mathcal{O}}$ est nulle}, et $D = J.E$ est donc un sous-fibré en droites de $E$ tel que $Q \,|\, D = 0$. Inversement, tout sous-fibré en droites de $E$ ayant cette propriété provient d'une section unique de $X'$ sur $X$. (Pour le vérifier, OPS $E = A$, et $Q$ la norme, et c'est alors immédiat.)

4. La notion de structure de similitude sur $E$, et la distinction de celles qui sont fidèles, str.\ fidèles, \add{non} dégénérées, « ne dépend que de $E$ modulo tensorisation par un module inv. », donc ne dépend que de l'algèbre d'Azumaya $M$ associée,
\[
  M = \underline{\mathrm{End}}(E),
\]
de rang 4 sur $E$. Si on part d'une telle $M$, la donnée d'une structure de similitude \uncertain{équivaut} \struck{relatif : un Module} à celle d'un sous-Module loc.\ \uncertain{monogène} \add{$\Delta$} de
\[
  M / \underline{\mathcal{O}}\, 1_M \overset{\mathrm{df}}{=} \Sigma(M),
\]
(qui joue le rôle du Module des formes quadratiques).

\page{28}
Le cas \struck{\ill{}} fidèle (\uncertain{resp.}\ str.\ fidèle) est celui où ce sous-Module est un Module inversible, resp.\ un sous-fibré en droites. Dans ce dernier cas, il \struck{\ill{}} y a un \struck{fibré} \add{rev.} quadratique associé de $X$, qui peut s'interpréter comme l'intersection de
\[
  \check{\mathbb{P}}(A) \hookrightarrow \check{\mathbb{P}}(M)
\]
(où $A \supset \underline{\mathcal{O}}.1_M$ est le sous-Anneau de $M$ image inverse de $\Delta$) avec la quadrique
\[
  Q \subset \check{\mathbb{P}}(M)
\]
définie par l'équation $\mathrm{Nr}_{M/\mathcal{O}}(x) = 0$, où $\mathrm{Nr}$ désigne la norme réduite. (C'est tautologique, la norme réduite induit sur $A$ celle de $A$.)

Quant au $A$-Module inv.\ \add{$E$}, i.e.\ le $\underline{\mathcal{O}}_{X'}$-Module inversible $\Lambda$, de la théorie précédente ? \struck{Il} Il est clair qu'il ne sera plus défini que modulo remplacement de $E$ par $E \otimes_{\underline{\mathcal{O}}_X} \Theta$, $\Theta$ inv.\ sur $X$, i.e.\ que $\Lambda$ sera défini modulo changement en
\[
  \Lambda \otimes p^{*}(\Theta),
\]
où $p \colon X' \to X$ le morphisme structural. Mais si $\sigma$ désigne l'inv.\ can.\ de $X'/X$, alors $\Lambda \otimes \sigma^{*}(\Lambda)^{-1}$ reste invariant, par changement de $\Lambda$ en $\Lambda \otimes p^{*}(\Theta)$. D'autre part,

\page{29}
on a
\[
  \mathrm{N}(\Lambda \sigma(\Lambda)^{-1}) \simeq \mathrm{N}(\Lambda)\bigl(\mathrm{N}(\sigma(\Lambda))\bigr)^{-1} \simeq \underline{\mathcal{O}} ,
\]
\uncertain{pour} ce dernier fait que
\[
  \mathrm{N}(\sigma(\Lambda)) \simeq \mathrm{N}(\Lambda) \quad \text{canoniquement.}
\]
Donc on s'attend : torseur \uncertain{sous} $A$-Module loc.\ libre, \uncertain{avec} restriction du groupe structural de $W(A)^{*}$ à $\mathcal{U}(A)$, \add{schéma en} groupes des sections locales de $A$ de norme 1.

Les $A$-Modules qui se proposent à la vue \uncertain{sont} \uncertain{évidemment} $M/A$, lequel est bien loc.\ libre de rang 1 sur $A$, si $A$ est un sous-fibré \struck{de $M$} \add{vect.\ de rang 2} (cas « str.\ fidèle »). \struck{Mais} Dans le cas où $M = \underline{\mathrm{End}}(E)$, il faudrait voir qu'on a un « iso.\ à scalaire près » canonique
\[
  E \simeq \underline{\mathrm{End}}(E)/A
\]
(iso.\ de $A$-modules, mais le « scalaire près » \uncertain{réfère} à une section de $\underline{\mathcal{O}}_X^{*}$…), \struck{on \uncertain{trouve} \uncertain{bien} \uncertain{nos} \ill{} \uncertain{schémas}}
\[
  \text{\struck{$\check{\mathbb{P}}(E) \simeq \mathbb{P}(\underline{\mathrm{End}}(E)^{+}/A)$}}
\]
\struck{on \uncertain{trouve} (\uncertain{comme}}
\[
  \text{\struck{$\mathbb{P}(E) \simeq \{$droites vect.\ $D$ dans $\underline{\mathrm{End}}(E)^{+}$ telles que $\mathrm{N}(D) = 0\}$,}}
\]
\[
  \text{\struck{$D \mapsto D^{\perp} \otimes D$),}}
\]
\marginal{où $\underline{\mathrm{End}}^{+}$ désigne le sous-Module des endom.\ de trace nulle}
\note{la fin de la page, depuis « on trouve », est barrée de traits obliques, la note marginale comprise ; un trait vertical en pointillés court le long de la marge gauche}

\page{30}
il s'agit de dire une iso.\ can.
\[
  \mathbb{P}(\underline{\mathrm{End}}(E)/A) \simeq C,
\]
\struck{\ill{}} \add{$C$ conique can.\ dans $\mathbb{P}(\underline{\mathrm{End}}^{+}(M))$}\note{le membre de droite, d'abord écrit puis biffé, est remplacé par $C$ en interligne}
(\emph{NB} en termes de $E$ \uncertain{cependant}, \uncertain{mais} en termes de $A \subset M$ seulement, un \struck{\ill{}} $C$ \uncertain{n'est} \uncertain{pas} en général \uncertain{iso.}\ à un $\mathbb{P}(\text{fibré vect.\ de rang } 2)$…)

Or on a sur $\underline{\mathrm{End}}(E)$ \emph{deux} structures de $A$-module, la gauche et la droite, échangées par l'antiinvolution canonique
\[
  u \mapsto \mathrm{Tr}(u) - u \qquad (\mathrm{Tr} = \text{trace réduite})
\]
et on a donc égalité des \uncertain{modules} \uncertain{déterminants}, \struck{\ill{}} $\underline{\mathrm{End}}(E)$ \struck{\ill{}} \add{$\underline{\mathrm{End}}(E)/A$} sur $A$, lesquels \uncertain{sont} \uncertain{égaux} can.\ :
\[
  \underline{\mathrm{End}}(E)/A \simeq \det\nolimits_A \underline{\mathrm{End}}(E) .
\]
Mais l'iso
\[
  \underline{\mathrm{End}}(E) \simeq E \otimes \check{E}
\]
transforme la structure de $A$-module gauche en celle définie par l'application $u_E \otimes \mathrm{id}_{\check{E}}$ dans, et la droite en $\mathrm{id}_E \otimes {}^{t}u_E$. Donc dans le cas de $A$-modules à gauche
\[
  \det\nolimits_A \underline{\mathrm{End}}(E) \simeq E^{\otimes_A 2} \otimes (\underbrace{\det\nolimits_h \check{E}}_{\check{L}})^{2}
\]
et à droite
\[
  \det\nolimits_A \underline{\mathrm{End}}(E) \simeq (\underbrace{\det\nolimits_h(E)}_{L})^{2} \otimes \check{E}^{\otimes_A 2},
\]
d'où un iso de $A$-modules
\[
  \underline{\mathrm{End}}(E)/A \simeq E^{\otimes_A 2} \otimes_h \check{L} \simeq \check{E}^{\otimes_A 2} \otimes_h L
\]
\note{la dernière formule est encadrée ; l'indice $h$ est lu tel quel, sans doute pour la base $\underline{\mathcal{O}}_X$}

\page{31}
\uncertain{En} \uncertain{dérive} que l'iso
\[
  \check{E}^{\otimes_A 2} \otimes_h \check{L} \simeq \check{E}^{\otimes_A 2} \otimes_h L
\]
i.e.
\[
  \check{E}^{\otimes_A 2} \simeq E^{\otimes_A 2} \otimes_h L^{\otimes_h (-2)}
\]
\note{la première formule est lue telle quelle ; le membre de gauche porte une surcharge, et la comparaison avec p.~30 suggère $E^{\otimes_A 2}$}
est \uncertain{déduit} de l'iso can.\ (qui respecte les structures de $A$-modules)
\[
  \check{E} \simeq E \otimes_h L^{\otimes_h (-1)}
\]
en \uncertain{élevant} \uncertain{au} \uncertain{carré} \uncertain{tensoriel} sur $A$. L'iso de $A$-modules cherché (mod \uncertain{scalaire} provenant de $\underline{\mathcal{O}}_X^{*}$)
\[
  E \simeq \underline{\mathrm{End}}(E)/A \qquad (\simeq E^{\otimes_A 2} \otimes_h \check{L})
\]
équivaut à \uncertain{la} \uncertain{donnée} d'une section de
\[
  E \otimes_h \check{L},
\]
\struck{base} de ce $A$-module mod.\ sections provenant de $\mathcal{O}^{*}$. Or il n'y a pas de canonique, visiblement, \struck{\ill{}} \add{i.e.\ il} n'y a pas de \uncertain{sous-fibrés} en droites canoniques $D \subset E$ \uncertain{associés} à la structure de similitude de départ sur $E$. Donc la supposition du début est erronée — $\underline{\mathrm{End}}(E)/A$ est, \uncertain{module} Module inversible provenant de la base, un Module inversible \emph{connu}, et c'est le \uncertain{seul} connu qu'il faudrait envisager à \uncertain{savoir}…

\page{32}
En rectifiant le tir, on trouve le $A$-module
\[
  E \otimes_A (\check{E})^{\otimes_A -1} \simeq M/\text{\struck{\ill{}}}\,M^{\natural}
\]
\note{le signe en exposant de $M$ ressemble à un bécarre ; il est transcrit $M^{\natural}$}
\uncertain{qui} \uncertain{est} \uncertain{inv.}\ sur $\underline{\mathcal{O}}$, où $M^{\natural}$ est défini localement sur $X$ par le \uncertain{choix} d'une base $(1, u)$ de $A$, via
\[
  \varphi_u \colon M \to M, \qquad \text{\struck{$\theta$}}\ \theta \mapsto u\theta - \theta\bar{u},
\]
\marginal{\emph{NB} $\varphi_u$ \uncertain{dérivation} de $M$ \ill{} $A$-\ill{} \ill{} \uncertain{Anneaux} \uncertain{modules} sur $M$}
\note{la note marginale, écrite en biais dans la marge gauche, est très difficile ; elle semble préciser la nature de $\varphi_u$}
(\emph{NB}\({}_2\) $\varphi_u$ dépend linéairement de $u$, et $\varphi_1 = 0$, donc \uncertain{on} $\varphi_u$ ne dépend que de la section image de $u$ \uncertain{dans} $\check{\Lambda} \simeq A/\underline{\mathcal{O}}$ \ill{} $M$ défini par $u$…), \struck{\ill{}} \uncertain{posant}
\[
  M^{\natural} = \mathrm{Im}\, \varphi_u, \quad \text{donc } M/M^{\natural} = \mathrm{Coker}\, \varphi_u .
\]
\uncertain{De} \uncertain{façon} plus intrinsèque, \ill{} suite exacte
\[
  M \otimes \check{\Lambda} \to M \to \Delta \to 0
\]
\note{la flèche de gauche a deux pointes ; le terme du milieu porte une rature}
\uncertain{visiblement} exacte de $A$-modules, pour la structure de $A$-Module \emph{gauche} de $M$). Il faudrait \uncertain{maintenant} définir sur $\Delta$ une forme quadratique \uncertain{non} \add{\uncertain{triviale}} \uncertain{canonique}, \struck{plus} précisément un iso.
\[
  \mathrm{N}_{A/\underline{\mathcal{O}}}(\Delta) \simeq \underline{\mathcal{O}} .
\]

\struck{\ill{}}\uncertain{Notons} sur $M$ l'autodualité donnée par
\[
  \langle \theta, \theta' \rangle \overset{\mathrm{def}}{=} \mathrm{Tr}\, \theta\bar{\theta}' = \mathrm{Tr}\, \theta'\bar{\theta}
\]
donc, si $\lambda_u$, $\gamma_u$ sont les opérations de la multiplication par $u$ à g.\ et à dr.
\[
  \lambda_u \theta = u.\theta, \qquad \gamma_u \theta = \theta u,
\]
on a
\[
  {}^{t}\lambda_u = \lambda_{\bar{u}}, \qquad {}^{t}\gamma_u = \gamma_{\bar{u}}
\]

\page{33}
donc la transposée de \struck{\ill{}}
\[
  \varphi_u \colon \theta \mapsto u\theta - \theta\bar{u}
\]
est
\[
  {}^{t}\varphi_u = \varphi_{\bar{u}} \colon \theta \mapsto \bar{u}\theta - \theta u ,
\]
\note{le signe devant le second $\varphi$ est peu net ; lu ${}^{t}\varphi_u = \varphi_{\bar{u}}$}
\uncertain{mais} :
\[
  \varphi_{\bar{u}} = -\varphi_u \quad \text{car } \bar{u} \equiv -u \bmod \underline{\mathcal{O}}.1 \text{ et } \varphi_{\mathcal{O}} = 0 .
\]
Ceci dit, on voit donc que \uncertain{dire} \uncertain{que} \emph{\uncertain{sur} $h$} \uncertain{de} $\mathrm{Coker}\, \varphi_u$ s'identifie à $\mathrm{Ker}\, {}^{t}\varphi_u$ $= \mathrm{Ker}\, \varphi_{\bar{u}}$ $= \mathrm{Ker}\, (-\varphi_u)$ $(= \mathrm{Ker}\, \varphi_u)$. Donc modulo Module inv.\ sur $\underline{\mathcal{O}}$, $\mathrm{Coker}\, \varphi_u$ s'identifie, en tant que $A$-Modules, au dual \add{sur $A$} de $\mathrm{Ker}\, \varphi_u$.

\struck{Soit}
\[
  B \subset M, \qquad B = \{\theta \in \Gamma M \mid \bar{u}\theta - \theta u = 0\}
\]
\marginal{pour une section locale de $A$ qui forme base de $A$ avec 1, donc pour \emph{toute} section de $A$ \uncertain{?}}
On a
\[
  B \cap M^{*} = \{\theta \in \Gamma M^{*} \mid \theta \text{ \uncertain{normalise} } A^{*} \text{ (ou $A$), et induit sur $A$ l'invol.\ can.\ } \sigma\}
\]
On voit donc que
\[
  B \cap M^{*} \text{ est un \add{pseudo-}torseur sous } \mathrm{Cent}_{M^{*}}(A) = A^{*} ;
\]
\note{l'égalité $\mathrm{Cent}_{M^{*}}(A) = A^{*}$ est écrite en exposant puis répétée en dessous}
\uncertain{Et} c'est un torseur \uncertain{ssi} $\exists$ section loc.\ \uncertain{partout}. Mais \struck{\ill{}} OPS $M = \mathbb{M}_2(\mathcal{O})$, et \struck{\ill{}} $A = \mathcal{O}[u]$, avec $u^{2} + bu + c = 0$, et \uncertain{comme}
\[
  \mathrm{Tr}\,\bar{u} = \mathrm{Tr}\, u = b, \qquad \mathrm{N}(\bar{u}) = \mathrm{N}(u) = c,
\]
\note{le signe de $\mathrm{Tr}\, u$ diffère de celui de p.~34 ($-b$) ; transcrit tel quel}
on sait que $\exists\, \theta \in M^{*}$ tel que $\struck{\ill{}}\ \mathrm{int}(\theta).u = \bar{u}$.

\page{34}
\emph{Dém.} OPS $A = \mathcal{O}[u]/(u^{2} + bu + c)$, et un automorphisme $\theta$ de la forme \uncertain{mentionnée}. \uncertain{Celui-ci} transforme $u$ en
\[
  u' = \bar{u} + \alpha.1_A \qquad \alpha \in \Gamma\underline{\mathcal{O}}
\]
(puisque $u' \equiv \bar{u} \equiv -u \bmod \underline{\mathcal{O}}.1_A$), avec
\[
  \left\{\begin{array}{l}
    \mathrm{Tr}(u') = \mathrm{Tr}(u) \ (= -b) \\
    \mathrm{N}(u') = \mathrm{N}(u) \ (= c)
  \end{array}\right.
\]
\marginal{la condition l'identité \uncertain{qui} \uncertain{l'inverse} \uncertain{soit} \ill{} \ill{} \ill{} \ill{} $\mathcal{U}$ \ill{} \uncertain{automorphismes} \ill{} $X_0$ \ill{}}
\note{note marginale écrite en biais dans la marge gauche, en grande partie illisible ; un trait la rattache au paragraphe suivant}
\struck{\ill{}} \uncertain{la} question correspondante \uncertain{bien} \uncertain{posée} : \uncertain{existe-t-il} $u'$ (\uncertain{sous} \uncertain{condition} \uncertain{sur} $A$)*. \struck{\ill{}} \uncertain{En} \uncertain{termes} de $\alpha$, les conditions s'écrivent
\[
  \left\{\begin{array}{l}
    2\alpha = 0 \\
    \alpha(\alpha - b) = 0
  \end{array}\right.
  \qquad (\text{car } \mathrm{Tr}\, \alpha.1_A = 2\alpha)
\]
\struck{\ill{}} \uncertain{Sur} les pts de $X - X_0$, \struck{\ill{}} \add{$2\alpha = 0$} implique $\alpha = 0$. \uncertain{Reste} : \uncertain{regarder} \uncertain{en} un point $s$ de $X_0$. \struck{\ill{} \ill{} que $\alpha_s \in \mathfrak{m}_s = $ rad $\underline{\mathcal{O}}_{X,s}$ — \uncertain{disons}, \ill{} le corps résiduel $k(s)$, \ill{} \ill{}}

\struck{\uncertain{dans} $\alpha(s)(\alpha(s) - b(s)) = 0$, mais par hypothèse $A(s)$ étale i.e.\ $b(s) \neq 0$, donc on doit avoir \uncertain{ou} \uncertain{bien} $\alpha(s) = 0$, ou bien $\alpha(s) = b(s)$};
\note{ce passage est barré de traits obliques ; un trait le relie à la suite, qui le remplace}
$\alpha_s$ serait inversible, et \uncertain{en} \add{la relation} \uncertain{aurait}, à cause de $\alpha(\alpha - b) = 0$, $\alpha_s - b_s = 0$ i.e.\ $\alpha_s = b_s$, \struck{\ill{}} i.e.\ $\alpha = b$ au voisinage de $s$, \struck{\ill{}} et comme $2\alpha = 0$ on \uncertain{aurait} $2b = 0$, donc comme $b_s$ est \emph{inversible} (car $A$ étale \uncertain{en} $s$, et $s \in X_0$), \struck{2.\ \ill{}} $1_X$ serait \uncertain{nul} au \uncertain{voisinage} de $s$ i.e.\ $X$ de car.\ 2.

\page{35}
\note{la symétrie $\sigma_e$ dont parle cette page est définie p.~36 ; les deux feuillets semblent intervertis, ou la page reprend un argument antérieur au lot}
\emph{\uncertain{Définition}} $\sigma_e$ s'appelle \emph{symétrie} p.r.\ à la \emph{droite} $D$, comme $\sigma$ est un automorphisme de $A$ pour la structure qu'il définit par le \uncertain{même}, $\sigma_e$ est un automorphisme de la structure quadratique de $E_x$ par $Q$, a fortiori un automorphisme de la structure de similitude.

\emph{Prop.} Supposons qu'en tout point de $X_0 = V(2.1_{\mathcal{O}_X})$, i.e.\ \uncertain{tout} pt de car.\ résid.\ 2, la structure de similitude \struck{\ill{}} sur $E$ soit non dég.\ (i.e.\ $A$ soit étale sur $\underline{\mathcal{O}}$ en \uncertain{ce} point). Alors pour \uncertain{toute} droite \struck{\ill{}} \uncertain{isotrope} $D \subset E$, $\sigma_D$ est l'\emph{unique} automorphisme $\theta$ de la structure de similitude $\Sigma$ (ou, ce qui revient au même, de la structure quadratique, quand $\Sigma$ est donnée par une forme quadratique $Q$) \struck{telle} qui induise $\mathrm{id}_{\uncertain{D}}$ sur $D$, $-\mathrm{id}$ sur $E/D$, et \add{qui} \uncertain{\emph{NB}} \add{soit} \struck{\ill{}} \add{l'identité au voisinage} \struck{\ill{}} \struck{$A_x$ \ill{} \ill{} $s \in S$}. (\uncertain{Cette} dernière condition \uncertain{n'a} \uncertain{lieu} \uncertain{qu'}en \uncertain{points} de $X - X_0$.)*
\struck{\ill{} \ill{} \ill{} \uncertain{identité} \ill{} \ill{} \ill{} \ill{}. \ill{} \ill{}}
\marginal{* \uncertain{Ceci} équivaut \uncertain{aussi} \uncertain{(}\uncertain{car} $\det \theta = -1$ \uncertain{et} $\theta \,|\, D = \mathrm{id}_D$\uncertain{)} \ill{} \ill{}}

\emph{Cor.} Soit $A$ algèbre quadratique sur $X$, supposons qu'en \uncertain{tout} $x \in X_0$, $A$ soit étale en $x$. Alors l'involution canonique est le seul automorphisme \add{\uncertain{d'alg.}} de $A$ qui induise $-\mathrm{id}$ sur $A/\underline{\mathcal{O}}.1_A$.

\page{36}
\struck{Alors}
\[
  f \colon A \to E, \qquad u \mapsto u.e
\]
est un iso respectant les structures quadratiques
\[
  Q(f(u)) = Q(u.e) = \mathrm{N}(u)\, \underbrace{Q(e)}_{1} = \mathrm{N}(u)
\]
\marginal{Plus généralement, $Q(u.x) = \mathrm{N}(u)\, Q(x)$ ($u \in \Gamma A$, $x \in \Gamma(E)$)}
\note{la note marginale est écrite en biais à gauche, reliée par un trait à $Q$ ; « Alors », écrit au-dessus, semble remplacé}
par quoi la structure d'alg.\ de $A$ est transformée \add{quadratique} en l'unique str.\ d'Alg.\ sur $\underline{\mathcal{O}}$, admettant $e$ pour section unité, et $Q$ comme fonction quadratique Norme. La symétrie cherchée $\sigma_D$ peut être définie, par transport de structure, comme déduite de l'inv.\ canonique $\sigma$ de $A$
\[
  \sigma \colon A \to A, \qquad \sigma(u) = \mathrm{Tr}(u) - u .
\]
Donc on aura, par définition (posant $\sigma_e = \sigma_{\underline{\mathcal{O}}e} = \sigma_D$)
\[
  \sigma_e(ue) = \bar{u}.e
\]
On \uncertain{va} donner une description équivalente, par
\[
  \sigma_e(x) = \varphi_Q(e, x)\,e - x
\]
où $\varphi_Q$ est la forme bil.\ \add{symétrique} associée à $Q$. \uncertain{Cette} formule \uncertain{se} \uncertain{fait} \uncertain{que} transformée \uncertain{celle} pour $A$, $\sigma(u) = \mathrm{Tr}(u) - u$, compte tenu que
\[
  \mathrm{Tr}(u) = \varphi_{\mathrm{N}}(u, 1) .
\]
Avec la notation du « produit scalaire » (qu'on ne \uncertain{divise} pas par 2), on aura $\varphi_Q(e, x) = (e.x)$, et on peut écrire
\[
  \sigma_e(x) = (e.x)\,e - x
\]
\note{le dernier membre est surchargé ; un « 2 » semble ajouté devant $(e.x)$, sans qu'on puisse dire s'il est maintenu}

\page{37}
\note{cette page et la suivante ne continuent pas p.~36 ; la fin de p.~37 (« telle que $Q(e) = 1$ ») est reprise en tête de p.~36, et p.~38 achève la démonstration de p.~34 : l'ordre des feuillets du fonds n'est pas celui de la rédaction}
\struck{J'ai dit} Le torseur en question provient d'un torseur $\Pi$ sous \struck{\ill{}} $\mathcal{U}_A \subset A^{*}$, sous-schéma en groupes des sections de norme 1, en posant
\[
  \Pi = \{\theta \in \Gamma\, \underbrace{\mathrm{S}M^{*}}_{\substack{\text{sections de}\\ \text{Norme } 1\\ (\mathrm{N}(\theta) = 1)}} \mid \theta \text{ \struck{\ill{}} normalise } A \text{ et y induit } u \mapsto \bar{u}\}
\]

\subsection*{4.2. Relation avec les symétries.}

Soit $(E, \Delta', Q)$, $Q \colon E \to \Delta'$, une structure de similitude str.\ fidèle. Soit $D$ un sous-fibré en droites de $E$, \struck{\ill{}} non isotrope i.e.\ $Q \,|\, D$ est non nul sur chaque fibre. Je vais définir un automorphisme
\[
  \sigma_D \colon E \xrightarrow{\sim} E
\]
de $E$ (\uncertain{pour} sa structure de similitude), induisant $\mathrm{id}_D$ sur $D$, $-\mathrm{id}$ sur $E/D$. \uncertain{Quitte} à tensoriser $E$ par $D^{-1}$, OPS $D$ muni d'une section base $e$, et profiter de celle-ci pour prendre une base $Q(e)$ dans $\Delta'$, donc OPS que $Q$ est \struck{défini par} une fonction quadratique scalaire
\[
  Q \colon E \to \underline{\mathcal{O}}
\]
telle que
\[
  Q(e) = 1 .
\]

\page{38}
\note{suite de la démonstration de p.~34}
au voisinage de $s$. Donc, au voisinage
\[
  u' = \underbrace{\bar{u}}_{-b-u} + b = \text{\struck{$(\mathrm{Tr}(u)1 - u)$}}\ -u = u
\]
et donc on aurait $\theta = \mathrm{id}$ au voisinage de $s$, contrairement à l'hypothèse.

Donc on conclut $\alpha_s \in \mathfrak{m}_s$, mais alors, comme $b_s \notin \mathfrak{m}_s$, \struck{\ill{}} on aurait $\alpha_s - b_s \notin \mathfrak{m}_s$ i.e.\ $\alpha_s - b_s \in \mathcal{O}_s^{*}$, d'où (comme $\alpha_s(\alpha_s - b_s) = 0$) $\alpha_s = 0$, donc $\alpha$ nul au voisinage de $s$, cqfd.

\emph{NB} En car.\ 2, dans le cas non étale, l'exemple typique est $A = \underline{\mathcal{O}}[u]/(u^{2})$. Les automorphismes de $A$ \add{(induisant $-\mathrm{id}$ sur $A/\underline{\mathcal{O}}.1$)} \struck{\uncertain{donnés}} \add{\uncertain{(i.e.\ $\mathrm{id}$ en car.\ 2)}} sont donnés par
\[
  \theta_{\lambda, \alpha} \colon u \mapsto \lambda u + \alpha \qquad \lambda \in \Gamma\, \mathcal{O}^{*},\ \alpha \in \Gamma\, \underline{\mathcal{O}}
\]
avec
\[
  \left\{\begin{array}{l}
    2\alpha = 0 \\
    \alpha^{2} = 0
  \end{array}\right.
\]
en car.\ 2, se réduit à :
\[
  \alpha^{2} = 0 .
\]
\uncertain{Les} \uncertain{automorphismes} qui induisent $-\mathrm{id}_{A/\mathcal{O}.1}$ sont ceux de la forme
\[
  u \mapsto -u + \alpha \qquad (\text{avec } 2\alpha = 0,\ \underline{\underline{\alpha^{2} = 0}}),
\]
donc en car.\ 2 \struck{\uncertain{ceux} de la forme} $u \mapsto u + \alpha$ avec $\alpha^{2} = 0$.
\marginal{\uncertain{Donc} \uncertain{l'involution} \ill{} \ill{} est \uncertain{la} \uncertain{seule} \ill{} qui \ill{} \uncertain{identité} \ill{} \ill{} \ill{} \ill{} \ill{} $2,0$ \ill{} \ill{} \ill{} $-1$ \ill{} \ill{} $\mathrm{id}$ \ill{}}
\note{note marginale écrite en biais dans la marge gauche, sur plusieurs lignes, en grande partie illisible}

Soit $(E, \Delta, Q)$ une str.\ de similitude str.\ fid., $D \subset E$ droite \uncertain{tot\textsuperscript{t}} non isotrope, d'où l'automorphisme $\sigma_D \in \Gamma\, \underline{\mathrm{End}}(E)^{*}$. Le fait que $\sigma_D$ soit un automorphisme de la structure de similitude signifie

\page{39}
signifie que
\[
  \underbrace{\mathrm{int}(\sigma_D)}_{\substack{\text{autom.\ de } \underline{\mathrm{End}}(E)\\ \theta \mapsto \sigma_D \theta \sigma_D^{-1}}} \text{ transforme } \underset{\underline{\mathrm{End}}(E)}{\overset{}{A}} \text{ en lui-même},
\]
\note{sous $A$ est écrit $\cap\ \underline{\mathrm{End}}(E)$ : $A$ est vu comme sous-Algèbre de $\underline{\mathrm{End}}(E)$}
et y \uncertain{est} \uncertain{l'inv.}, plus précisément
\[
  \boxed{\sigma_D u \sigma_D^{-1} = \bar{u}} \qquad \text{pour } u \in \Gamma A .
\]
Pour le prouver, OPS $E = A$, \add{$\Delta = \underline{\mathcal{O}}$, $Q = \mathrm{N}_{A/\underline{\mathcal{O}}}$}, $D = \underline{\mathcal{O}}.1$, et on est ramené à prouver que
\[
  \sigma \mu_u \sigma^{-1} = \mu_{\bar{u}}
\]
où $\mu_u \in \Gamma\, \underline{\mathrm{End}}(A)$ est défini par $\mu_u x = ux$. La formule équivaut alors :
\[
  \overline{u\bar{x}} = \bar{u}\, x \qquad \text{OK.}
\]

\emph{Proposition} $(E, \Delta, Q)$, ou encore $(E, A \subset \underline{\mathrm{End}}(E))$ similitude str.\ fid. Soit
\[
  \Pi = W(\underline{\mathrm{End}}(E))^{*}
\]
\[
  \overset{\shortparallel}{\{}\theta \text{ sections locales de } \underline{\mathrm{End}}(E) \text{ dans } \mathrm{Sch}_{/X} \mid \theta \text{ normalise } A, \text{ induit } \sigma \text{ sur } A \text{ et } \det\theta = -1\},
\]
\struck{Alors} \add{et soit} \struck{\ill{}}
\[
  M = W(\underline{\mathrm{End}}(E))
\]
\[
  \overset{\shortparallel}{\{}\theta \mid \theta u = \bar{u}\theta \text{ pour toute section } u \text{ de } A\}
\]
(ou pour une section qui engendre $A/\underline{\mathcal{O}}.1_A$ …)
\note{dans la première accolade, l'étoile de $W(\underline{\mathrm{End}}(E))^{*}$ porte un renvoi qui ne semble pas résolu sur la page}

\page{40}
Alors

a) $M$ est un sous-fibré vectoriel de rang 2 de $\underline{\mathrm{End}}(E)$, qui est un sous-$A$-Module à gauche et à droite. Les deux structures \uncertain{gauche} et droite sont reliées par
\[
  u\theta = \theta\bar{u}, \qquad \text{pour } \theta \in \Gamma M,\ u \in \Gamma A
\]
(elles se \uncertain{déterminent} \uncertain{donc} \uncertain{mutuellement} …). Le $A$-Module $M$ (à gauche \add{ou à droite}) est \struck{loc.\ libre} inv.\ (\uncertain{comme} \ill{}).

b) $\Pi$ est formé des sections $\theta$ de $M$ telles que
\[
  \mathrm{N}(\theta) = -1 .
\]
Il forme un torseur* sous $\mathrm{S}W(A) \subset W(A)^{*}$, le \struck{groupe} schéma en groupes des $x \in \Gamma A$ telles que $\mathrm{N}(x) = 1$, par l'opération de multiplication : \uncertain{g.}\ (\uncertain{ou} \uncertain{à} \uncertain{droite}, par \uncertain{celle} \uncertain{à} droite).
\marginal{* \emph{NB} \uncertain{Ce} \uncertain{torseur} \uncertain{fait} \uncertain{que} \ill{} \uncertain{triviale} (\uncertain{Zar.})}

\struck{\emph{Dém.} OPS $E = A = \underline{\mathcal{O}}[u]/(u^{2} + bu + c)$. Alors $W$ est défini comme le \ill{} de \ill{} $\underline{\mathrm{End}}$ \ill{}}
\note{les trois lignes de cette démonstration sont barrées de traits obliques}

c) Considérons
\[
  \underset{\overset{\cup}{X'}}{P} = \mathbb{P}(E), \qquad P^{*} = P - \underset{\mathrm{Spec}(A)}{\overset{\shortparallel}{X'}}
\]
le fibré des droites \uncertain{tot\textsuperscript{t}} non isotropes de $A$. $W(A)^{*}$ opère sur $P^{*}$ \uncertain{grâce} \uncertain{à} \uncertain{son} opération sur $E$, \uncertain{où} $u \in \Gamma\, W(A)^{*}$ définit la transformation $D \mapsto u.D$ (\emph{NB} $\mathbb{G}_m = W(\underline{\mathcal{O}})^{*} \subset A^{*}$ opère trivialement donc c'est $W(A)^{*}/\mathbb{G}_m$ qui opère…), \uncertain{donc} \struck{application} une application \uncertain{équivariante} canonique
\[
  \begin{array}{rcl}
    P^{*} & \xrightarrow{f} & \Pi \\
    D & \longmapsto & \sigma_D
  \end{array}
\]
\note{l'argument se poursuit au-delà du lot}

\end{document}
